\bigskip\Needspace{10\baselineskip}
\pdfbookmark[1]{Further Report. Ceres Ferdinandea. April 1802}{ceres180204}
\begin{center}
{\small FURTHER REPORT ON THE NEW MAJOR PLANET.\par}
CERES FERDINANDEA.\par
\medskip\hrule\medskip
Monatliche Correspondenz Vol. V. pp.\ 379—404. April 1802.\par
\medskip\hrule
\end{center}\medskip
—\quad —\quad —\quad \textsc{Gauss}'s elliptic orbit VII for the \textit{new planet} [p.\ 207 of this volume] continues to agree closely with our ongoing observations of this planetary body, as the following calculation and comparison by Dr.\ \textsc{Gauss} amply demonstrate.
\begin{center}
{\small\setlength{\tabcolsep}{5pt}
\begin{tabular}{r r|r|r|r|r}
\multicolumn{2}{c|}{\shortstack{Seeberg\\1802}}&\shortstack{Calculated\\right ascension\\$\CeresSymbol$}&Difference&\shortstack{Calculated\\northern\\declination}&Difference\\\hline
February&19&$187^0\ 58'\ 23''8$&$-4''1$&$14^0\ 20'\ 24''0$&$+21''1$\\
—&26&$187\quad7\ 34.4$&$+1.0$&$15\quad9\ 20.1$&$+25.8$\\
—&27&$186\ 58\ 53.2$&$+5.8$&$15\ 16\ 20.2$&$+26.3$\\
—&28&$186\ 49\ 51.9$&$+3.7$&$15\ 23\ 19.1$&$+18.1$\\
March&1&$186\ 40\ 31.4$&$+3.5$&$15\ 30\ 16.0$&$+34.9$\\
—&2&$186\ 30\ 52.2$&$+3.3$&$15\ 37\ 10.6$&$\ .\ .\ .$\\
—&3&$186\ 20\ 55.1$&$+2.6$&$15\ 44\quad2.0$&$+19.7$
\end{tabular}}
\end{center}
The right ascensions therefore remain in excellent agreement, and the error in declination seems to have changed little, if at all. With the observation of 1 March left out, its mean over five observations would be $+22''$. —\quad —\quad —

Oberamtmann Dr.\ \textsc{Schröter} continues to observe the diameter of \textit{Ceres} regularly. Here are some of his observations, which Dr.\ \textsc{Gauss} calculated and was kind enough to share with us.
\begin{center}
{\small\setlength{\tabcolsep}{4pt}
\begin{tabular}{r l r l r}
\multicolumn{3}{c}{Apparent diameter}&\multicolumn{2}{c}{\shortstack{Calculated diameter\\at distance 1.}}\\[5pt]
1802 January 10&$\ .\ .\ .\ .$&$2''500$&$\ .\ .\ .\ .\ .\ .$&$5''16$\\
25&$\ .\ .\ .\ .$&$2.514$&$\ .\ .\ .\ .\ \raisebox{.5ex}{.}\ .$&$4.77$\\
26&$\ .\ .\ .\ .$&$2.687$&$\ .\ .\ .\ .\ .\ .$&$5.07$\\
28&$\ .\ .\ .\ .$&$2.774$&$\ .\ .\ .\ .\ .\ .$&$5.18$\\
31&$\ .\ .\ .\ .$&$2.930$&$\ .\ .\ .\ .\ .\ .$&$5.38$\\
February 5&$\ .\ .\ .\ .$&$3.120$&$\ .\ .\ .\ .\ .\ .$&$5.59$\\
10&$\ .\ .\ .\ .$&$3.342$&$\ .\ .\ .\ .\ .\ .$&$5.90$
\end{tabular}}
\end{center}
\SourceLeaf{215}%

The mean would be $5''29=0.308$ of the Earth, or 529 geographical miles, including the. atmospheric haze. According to \textsc{Schröter}'s measurement of the sharply bounded disk, its diameter would be only $3''44$ at distance $=1$, so almost exactly $\tfrac15$ of the Earth's diameter, or considerably smaller than our Moon. No wonder, then, that such a little planet remained hidden from us for so long!

Since \textsc{Gauss}'s ellipse VII for the orbit of \textit{Ceres} still represents the observations of this planet fairly accurately, we are giving our astronomical readers, while there is still time, a continuation of the \textit{ephemeris} of its motion for the next three months. Dr.\ \textsc{Gauss} very kindly undertook the calculation. It is hard to believe that \textit{Ceres} will remain visible until then. Already, on 27 and 28 March, we are noticing such a rapid change in the light of this planet that we are inclined to attribute it to causes quite different from the state of our atmosphere or the planet's distance from the Earth. Still, for astronomers equipped with equatorial and parallactic instruments who wish to leave nothing untried, this will make the increasingly difficult search for Ceres considerably easier.
\begin{center}
Positions of \textit{Ceres} in April, May and June 1802 at midnight, mean Seeberg time.\par\medskip
{\fontsize{9}{11}\selectfont\setlength{\tabcolsep}{2pt}
\begin{tabular}{r r|r|r|r@{\qquad}r r|r|r|r}
\multicolumn{2}{c|}{}&\shortstack{R.A. $\CeresSymbol$\\in degrees}&\shortstack{Northern\\decl.}&\shortstack{R.A. $\CeresSymbol$\\in time}&\multicolumn{2}{c|}{}&\shortstack{R.A. $\CeresSymbol$\\in degrees}&\shortstack{Northern\\decl.}&\shortstack{R.A. $\CeresSymbol$\\in time}\\\hline
April&21&$177^0\ 14'$&$17^0\ 52'$&$11^{\mathrm u}\ 48^{\mathrm m}\ 55^{\mathrm s}$&May&27&$177^0\ 14'$&$14^0\ 13'$&$11^{\mathrm u}\ 48^{\mathrm m}\ 57^{\mathrm s}$\\
&24&$176\ 56$&$17\ 42$&$11\ 47\ 43$&&30&$177\ 34$&$13\ 47$&$11\ 50\ 16$\\
&27&$176\ 41$&$17\ 30$&$11\ 46\ 44$&June&2&$177\ 57$&$13\ 21$&$11\ 51\ 47$\\
&30&$176\ 29$&$17\ 17$&$11\ 45\ 58$&&5&$178\ 22$&$12\ 53$&$11\ 53\ 27$\\
May&3&$176\ 21$&$17\quad2$&$11\ 45\ 25$&&8&$178\ 49$&$12\ 25$&$11\ 55\ 17$\\
&6&$176\ 17$&$16\ 45$&$11\ 45\quad6$&&11&$179\ 19$&$11\ 56$&$11\ 57\ 16$\\
&9&$176\ 15$&$16\ 27$&$11\ 45\quad2$&&14&$179\ 51$&$11\ 27$&$11\ 59\ 24$\\
&12&$176\ 17$&$16\quad8$&$11\ 45\ 10$&&17&$180\ 25$&$10\ 57$&$12\quad1\ 40$\\
&15&$176\ 23$&$15\ 47$&$11\ 45\ 27$&&20&$181\quad1$&$10\ 27$&$12\quad4\quad4$\\
&18&$176\ 31$&$15\ 25$&$11\ 46\quad5$&&23&$181\ 39$&$9\ 56$&$12\quad6\ 35$\\
&21&$176\ 43$&$15\quad2$&$11\ 46\ 50$&&26&$182\ 18$&$9\ 25$&$12\quad9\ 13$\\
&24&$176\ 57$&$14\ 38$&$11\ 47\ 48$&&29&$183\quad0$&$8\ 54$&$12\ 11\ 59$
\end{tabular}}
\end{center}

\textit{Ceres} comes to a standstill and resumes direct motion on 2 May in longitude at $5^{\mathrm z}\ 19^0\ 45'$, and on 9 May in R.A. at $176^0\ 15'$. This time the retrograde motion amounts to $13^0\ 10'$ in longitude and lasts 92 days. In R.A. it amounts to $12^0\ 27'$ and lasts 93 days. Dr.\ \textsc{Gauss} calculated the relative brightness of the planet as follows, disregarding the magnitude of the phase.
\begin{center}
{\small\setlength{\tabcolsep}{5pt}
\begin{tabular}{r l r|r|r|r}
\multicolumn{3}{c|}{}&\multicolumn{2}{c|}{Distance from}&Brightness\\
\multicolumn{3}{c|}{}&$\odot$&$\EarthSymbol$&\\\hline
1801&December&7&$2.5453$&$2.4892$&$0.4227$\\
1802&March&16&$2.5706$&$1.6024$&$1.0000$\\
&May&6&$2.5961$&$1.9039$&$0.6945$\\
&June&29&$2.6305$&$2.5735$&$0.3703$
\end{tabular}}
\end{center}
\begin{flushright}[\textsc{von Zach}.]\end{flushright}
